\section{Literature Review}
\label{sec:literature}

\subsection{Ecological Resilience Measurement}

Ecological resilience, defined as the capacity of an ecosystem to
absorb disturbance and reorganise while retaining essentially the same
functions and feedbacks, has become a central concept in environmental
management \citep{ipcc2023cities}. Quantifying resilience, however,
remains methodologically contested. The Remote Sensing Ecological Index
(RSEI) integrates greenness, wetness, heat, and dryness components
derived from satellite imagery into a single composite score, and has
been widely applied to assess urban ecological quality across Chinese
cities \citep{chen2023evaluation,liu2024analysis,huang2024assessment}.
Recent advances have extended RSEI with deep-learning feature
extraction \citep{gong2025beyond} and water-body corrections
\citep{zhou2025dynamic}, while \citet{zhang2025urban} propose an
improved RSEI that addresses the computational challenges of its land
surface temperature component. \citet{zhang2025constructing} demonstrate
that RSEI can be integrated with circuit theory to construct ecological
security patterns, bridging the gap between quality assessment and
spatial planning.

A limitation of RSEI and similar composite indices is that they measure
ecological \emph{state} rather than ecological \emph{resilience}---the
dynamic response to disturbance. \citet{smith2023global} address this by
quantifying vegetation resilience through the temporal autocorrelation
of NDVI, finding that resilience is linked to water availability and
variability globally. \citet{li2023widespread} show that spring
phenology significantly affects drought recovery times across the
Northern Hemisphere, underscoring the importance of measuring recovery
as a distinct dimension of resilience. These studies motivate our
approach of decomposing ecological resilience into resistance (the
degree to which NDVI is maintained during an extreme event) and
recovery (the speed of NDVI restoration afterwards), rather than
relying on a static composite index alone.

\subsection{Climate Adaptation Policy Evaluation in China}

China has implemented a series of pilot programmes to promote urban
climate adaptation. The sponge city programme, launched in 2015--2016,
focused on stormwater management through permeable pavements, green
roofs, and bioswales \citep{han2023chinas}. \citet{wang2024sponge}
evaluate the sponge city pilot's effect on ecological livability,
finding modest improvements in green coverage. The broader
climate-adaptive city pilot programme, launched in 2017, extends
beyond stormwater to encompass heat-island mitigation, ecological
corridor construction, and comprehensive climate risk assessment
\citep{ipcc2023cities}. \citet{wang2024urban} provide an overview of
nature-based solutions in Chinese urban planning, including the
National Forest City programme, while \citet{lu2026does} examine the
effect of culture-tourism integration policy on urban ecological
resilience using multi-period DID.

A growing literature applies quasi-experimental methods to evaluate
these policies. \citet{gao2025renewable} use advanced data analysis
techniques to assess renewable energy policies' impact on greenhouse
gas emissions. \citet{yuan2024double} employ a DML model to measure the
effect of the ``Made in China 2025'' strategy on green economic growth,
demonstrating the applicability of machine-learning causal inference
to Chinese policy evaluation. \citet{qian2025impact} apply DML to
evaluate the zero-waste city pilot's effect on corporate green
transformation, while \citet{yu2025carbon} combine staggered DID with
DML to estimate the carbon mitigation effect of service trade
innovation. \citet{zhang2025does} use DML to study the synergistic
effect of digital infrastructure on pollution and carbon reduction.
These studies collectively demonstrate the potential of DML for policy
evaluation, but none has applied it specifically to the climate-adaptive
city pilot's ecological outcomes.

\subsection{Double Machine Learning for Causal Inference}

Double/debiased machine learning (DML) \citep{chernozhukov2022automatic}
provides a framework for estimating low-dimensional causal parameters in
the presence of high-dimensional or nonparametric nuisance functions.
The key innovation is the combination of Neyman orthogonality---which
insulates the target parameter from first-order bias in nuisance
estimation---with cross-fitting, which prevents overfitting by
separating the samples used for nuisance estimation and target
parameter estimation \citep{ahrens2024ddml,bia2023double}. This approach
has several advantages over conventional parametric DID. First, it
allows flexible machine-learning algorithms to model the relationship
between confounders and both the treatment and outcome, without
imposing linearity or additivity. Second, the orthogonality condition
ensures $\sqrt{n}$-consistency and asymptotic normality even when the
nuisance functions converge at slower rates. Third, it naturally
accommodates multiple ML learners, enabling algorithm comparison as a
form of specification robustness \citep{byrnes2025causal}.

The DID--DML hybrid we employ extends the standard DML partially linear
model by interacting the treatment indicator with a continuous measure
of climate shock intensity. This specification captures the policy's
differential effect under varying levels of climate stress, which is
the theoretical mechanism through which climate-adaptive infrastructure
is expected to operate: the policy's ecological benefits should be most
pronounced when cities face more severe extreme weather events.
\citet{callaway2024did} formalise DID with continuous treatments, while
\citet{borusyak2024revisiting} develop robust event-study estimators
for staggered designs. Our approach, by focusing on a single treatment
cohort (2017), avoids the negative-weight problems that plague staggered
TWFE estimators \citep{chaisemartin2023twoway,roth2024interpreting}.
